The overall score for the competition is the sum of all the challenge rounds performed by the Cyber Reasoning System (CRS).
Each round has a fixed duration $CT$ (which is usually between 6 hours and 12 hours), and each event is marked with the time remaining in the duration of the challenge, $RT$.

The score for a challenge round (CS) is determined by various factors:
\begin{equation}
CS = AM * (VDS + PRS + SAS + BDL),
\end{equation}
where $AM$ is the Accuracy Multiplier, $VDS$ is the Vulnerability Discovery Score, $PRS$ is the Program Repair Score, $SAS$ is the SARIF Assessment Score, and $BDL$ is the Bundle Score.

The Accuracy Multiplier represents and incentivizes the accuracy of the CRS. The Accuracy Multiplier is computed as follows:
\begin{equation}
AM = 1 - (1 - r)^4,
\end{equation}
where $r$ is the ratio of accurate submissions to all submissions:
\begin{equation}
r = (accurate)/(accurate + inaccurate).
\end{equation}

The $VDS$ represents the CRS' ability to find vulnerabilities and produce inputs that trigger them (PoVs), and it is calculated as the sum of all distinct PoV values:
\begin{equation}
VDS = \sum value_{PoV}, \forall PoV \in ScoringPoVs.
\end{equation}
The set of $ScoringPoVs$ is the set of unique vulnerabilities identified. 
If a PoV does not trigger a vulnerability, its value is 0.
If a submitted PoV is considered a duplicate of a previously submitted PoV, then the previous PoV is ignored, and the new PoV is considered as the entry for the associated vulnerability.
This is important as the submission time affects the overall score.
\begin{equation}
value_{PoV} = 2 * (0.5 + RT/(2*CT)).    
\end{equation}
Note that if a PoV is overwritten by a subsequent PoV, this does not affect the accuracy multiplier.

The $PRS$ represents the ability of the CRS to generate effective patches.
The score is calculated as the sum of all the scored patches:
\begin{equation}
PRS = \sum value_{patch}, \forall patch \in ScoringPatches.
\end{equation}

A patch that does not remediate any challenge vulnerability,  fails to build, or does not pass the functionality tests has a value of 0.
Usually, the evaluation of a patch by the organizers takes around 15 minutes.
The functionality tests for a patch are sometimes included with the challenge, but not always.
If a patch can be applied, remediates one or more challenge vulnerabilities, and passes the functionality tests, its value is the following.
\begin{equation}
value_{patch} = 6 * (0.5 + RT/(2*CT)).
\end{equation}
That is, similar to the case of PoVs, the value of a patch diminishes when it is submitted later in the round.

For a patch to be evaluated in the scoring process, it has to solve all variant collected PoVs created during the challenge round. 
This means that a patch submitted by system A could be invalidated because it does not block the PoV submitted by system B for the same vulnerability.
Note that if a patch is invalidated, it will contribute (negatively) to the Accuracy Multiplier.

The SARIF Assessment Score ($SAS$) measures the ability of the CRS to validate a SARIF report.
\begin{equation}
SAS = \sum value_{assessment}, \forall assessment \in ScoringAssessments.
\end{equation}

If a SARIF assessment is incorrect, then the corresponding value is 0.
The time penalty for the assessment is determined with respect to the remaining challenge time when the SARIF report was made public, which is $SRT$. 
If a SARIF assessment is correct, its value is the following:
\begin{equation}
value_{assessment} = 1 * (0.5 + RT/(2*SRT)).
\end{equation}

If multiple assessments are submitted for the same report, only the lastl be scored, and the previous ones will negatively affect the accuracy.

The Bundle Score ($BDL$) represents the ability of the CRS to associate PoVs, patches, and SARIF reports.
\begin{equation}
BDL = \sum value_{bundle}, \forall bundle \in ScoringBundles,
\end{equation}
where $ScoringBundles$ are bundles that contain two or more PoVs, patches, or SARIF reports.

The value of the bundle is computed as follows:
\begin{equation}
value_{bundle} = 0.5 * (value_{PoV} + value_{patch}) + b, 
\end{equation}
if both the PoV and patch are supplied; otherwise:
\begin{equation}
value_{bundle} = b,
\end{equation}
where b takes the following values:
\begin{description}
    \item[0:] if no SARIF ID is included in the bundle, or the bundle contains incorrect pairings;

    \item[1:] if the bundle has the correct pairing for the SARIF ID and the PoV but there is no patch included;

    \item[2:] if the bundle has the correct pairing for the SARIF ID and the patch but has no PoV;

    \item[3:] if the bundle contains correct pairings for SARIF ID, PoV, and patch.
\end{description}

The definition of a correct pairing is as follows:
\begin{itemize}
\item For a PoV and patch pairing the patch must remediate the PoV;
\item For a PoV and SARIF ID pairing, the PoV must crash the vulnerability defined in the SARIF report;
\item For a patch and SARIF ID pairing, the patch must fix the vulnerability described in the SARIF report;
\item For a PoV, patch, and SARIF ID pairing all of the above conditions must hold.
\end{itemize}
At the end of the challenge round, there can be only one bundle submitted per SARIF ID or challenge vulnerability.

Most importantly, if a bundle has an incorrect pairing, its value is detracted from the overall score, highlighting the fact that bundles represent high-risk/high-reward components of the score.
